% ============================================
% Johns Hopkins Letter Template
% Uses the built‐in 'letter' class
% ============================================

\documentclass[11pt,letterpaper]{letter}

\usepackage{graphicx}
\usepackage{eso-pic}
\usepackage{geometry}
\usepackage{microtype}
\usepackage[utf8]{inputenc}
\usepackage[hidelinks]{hyperref}

% ----- Page Setup -----
\geometry{left=25mm,right=25mm,top=25mm,bottom=25mm}

% ----- Logo Placement (foreground, first page only) -----
\newcommand{\PutJHULogoFG}[1]{%
  \AddToShipoutPictureFG*{%
    \AtPageUpperLeft{%
      \hspace{10mm}%
      \raisebox{-38mm}{%
        \includegraphics[width=6cm]{#1}%
      }%
    }%
  }%
}

% ----- Sender Information -----
\signature{Decory Edwards}
\address{Department of Economics \\ Johns Hopkins University \\ Baltimore, MD 21218 \\ \texttt{dedwar65@jh.edu}}

% ----- Recipient Info -----
\begin{document}
\PutJHULogoFG{Images/jhulogo.png} % show logo only on first page

\begin{letter}{Hiring Committee \\
Compass Lexecon \\
555 12th Street NW, Suite 501 \\
Washington, DC 20004}

\opening{Dear Hiring Committee,}

I am writing to express my interest in the Senior Economist position at Compass Lexecon. I am a Ph.D.\ candidate in the Department of Economics at Johns Hopkins University (advisors: Christopher D. Carroll, Nicholas W.\ Papageorge, and Stelios Fourakis), and I expect to receive my degree in May 2026. My primary areas of specialization are macroeconomic modeling, wealth and income dynamics, and computational economics.

My research agenda focuses on the ability of macroeconomic models to generate empirically realistic distributions of wealth and income over the life cycle. A key driver of that work is modeling heterogeneity in the rate of return on assets, and understanding how return dispersion contributes to portfolio accumulation across households. These themes align strongly with Compass Lexecon’s consulting practice, where rigorous economic and financial analysis supports legal, regulatory, and strategic decision-making.

In my job-market paper, I calibrate a life-cycle heterogeneous-agent model to match wealth moments from the Survey of Consumer Finances by allowing households to differ in their portfolio returns. The model helps explain observed wealth concentration and return dispersion. In a related project using the Health and Retirement Study, I examine how trust influences both income growth and asset returns — finding a hump-shaped relationship for returns, and characterizing the distribution of the persistent component of returns to net worth. These experiences have equipped me to: (i) frame research strategies and design modeling approaches; (ii) execute econometric and simulation analysis with large microdata sets; and (iii) translate results into clear, rigorous written deliverables for academic and practitioner audiences.

I am particularly drawn to Compass Lexecon’s emphasis on applying theory and quantitative methods to real-world cases — including damages estimation, merger analysis, regulatory matters and valuation issues. My quantitative toolkit (Python, R, Stata), my experience with heterogeneity and distributional outcomes, and my ability to distill complex models into actionable insights make me well-suited to contribute to high-stakes consulting engagements and client communication.

I believe I would be a strong fit at Compass Lexecon because:
\begin{enumerate}
    \item My expertise in simulation-based modeling and empirical validation aligns with analyzing complex datasets and building models used in litigation and strategic consulting.
    \item My research on heterogeneity, returns and wealth dynamics maps directly to consulting questions about valuation, financial behavior, and economic harm.
    \item I bring strong communication skills and experience working across teams, which are essential for delivering concise summaries, client-facing presentations, and collaborating in a consulting environment.
\end{enumerate}

I am excited about the prospect of joining Compass Lexecon’s team and contributing to rigorous consulting work that bridges academia and practice. Thank you for considering my application.

\closing{Sincerely,}

\end{letter}
\end{document}
